\ifdefined\gkver\else\def\gkver{student}\fi
\documentclass[\gkver]{gaokaozhenti}
\begin{document}

\chapter{2021年上海}

\begin{enumerate}
\item
%% number: 1
%% paperName: 2021年上海高考第 1 题 3 分
%% typeId: 1
%% type: 单选题
%% chapter: 原子和原子核
%% point: 天然放射性
%% method: 无
%% score: 3
%% degree: 624
%% duplicateId: 0
%% body:
$\beta$ 粒子\xzanswer{B}
\fourchoices[answer=B]
{是原子的外层电子}
{来自中子的转化}
{是原子核内的核子}
{来自质子的转化}
%% answer: B
%% memo:
\memoanswer{
$\beta$ 粒子是电子，$\beta$ 衰变的本质是原子核中的 1 个中子转变为 1 个质子和 1 个电子，对应的核反应方程为：$\ce{^{1}_{0}n -> ^{1}_{1}H + ^{0}_{-1}e}$。正确选项为 B。
}

\item
%% number: 2
%% paperName: 2021年上海高考第 2 题 3 分
%% typeId: 1
%% type: 单选题
%% chapter: 原子和原子核
%% point: 原子核式结构
%% method: 无
%% score: 3
%% degree: 860
%% duplicateId: 0
%% body:
在卢瑟福 $\alpha$ 粒子散射实验中，$\alpha$ 粒子发生大角度偏转\xzanswer{D}
\fourchoices[answer=D]
{主要是由于受到电子的碰撞}
{是由于受到原子核的库仑引力}
{表明原子核由质子和中子组成}
{反映了原子内存在带正电的核}
%% answer: D
%% memo:
\memoanswer{
电子的质量很小，与 $\alpha$ 粒子碰撞后几乎不会影响它的轨迹，选项 A 错误；

$\alpha$ 粒子是氦原子核，带正电，与同样带正电的原子核间存在相互排斥的库仑力，正是这个力使 $\alpha$ 粒子发生大角度偏转。选项 D 正确，B、C 错误。
}

\item
%% number: 3
%% paperName: 2021年上海高考第 3 题 3 分
%% typeId: 1
%% type: 单选题
%% chapter: 机械波
%% point: 波长频率波速
%% method: 无
%% score: 3
%% degree: 801
%% duplicateId: 0
%% body:
位于弹性绳中点 O 的波源振动频率随时间逐渐增大，经一段时间后弹性绳的形状可能为\xzanswer{A}
\fourchoices[answer=A, ispicture=true, w=6cm]
{\includesvg{figs/03a}}
{\includesvg{figs/03b}}
{\includesvg{figs/03c}}
{\includesvg{figs/03d}}
%% answer: A
%% memo:
\memoanswer{
O 为波源，它在绳上会产生两列分别向左、右传递的波，选项 B、D 错误；

波速由介质决定，在同一根弹性绳上这两列波波速相同，由 $\lambda=\frac{v}{f}$ 可知，随着频率的增大，波长会随之变小。正确选项为 A。
}

\item
%% number: 4
%% paperName: 2021年上海高考第 4 题 3 分
%% typeId: 1
%% type: 单选题
%% chapter: 力物体的平衡
%% point: 受力分析
%% method: 无
%% score: 3
%% degree: 552
%% duplicateId: 0
%% body:
我国“天问一号”火星探测器计划依靠发动机推力悬停于火星表面附近，以寻找合适的着陆点。当探测器在空中水平匀速移动时，发动机的喷气方向（不计火星大气阻力）\xzanswer{C}
\fourchoices[answer=C]
{竖直向上}
{斜向上}
{竖直向下}
{斜向下}
%% answer: C
%% memo:
\memoanswer{
探测器在水平方向上匀速移动，处于受力平衡状态，探测器推力的方向应竖直向上，则喷气方向为竖直向下。正确选项为 C。
}

\item
%% number: 5
%% paperName: 2021年上海高考第 5 题 3 分
%% typeId: 1
%% type: 单选题
%% chapter: 气体性质
%% point: 分子动理论
%% method: 无
%% score: 3
%% degree: 901
%% duplicateId: 0
%% body:
某理想气体在 $T_{1}$、$T_{2}$ 两个不同温度下的分子速率分布曲线如图所示。图中 $f(v)$ 表示单位速率区间内的分子数占总分子数的百分比随分子速率 $v$ 的变化关系，两曲线与横轴所围面积分别为 $S_{1}$、$S_{2}$，则\xzanswer{B}
\onepicture[h=3.6cm, align=right]{\includesvg{figs/05}}
\fourchoices[answer=B]
{$T_{1} > T_{2}$}
{$T_{1} < T_{2}$}
{$S_{1} > S_{2}$}
{$S_{1} < S_{2}$}
%% answer: B
%% memo:
\memoanswer{
温度越高，分子无规则运动越剧烈，或者说更加杂乱无章，具有最大比例的速率区间会向高温段偏移，所以 $T_{1}<T_{2}$；

此图像的面积的物理意义为所有速率分子所占概率之和，为 100\% 或者说成 1，$S_{1}=S_{2}$。

正确选项为 B。
}

\item
%% number: 6
%% paperName: 2021年上海高考第 6 题 3 分
%% typeId: 1
%% type: 单选题
%% chapter: 直线运动
%% point: 运动图象
%% method: 无
%% score: 3
%% degree: 870
%% duplicateId: 0
%% body:
如图，一作直线运动的质点在 $0<t<4\Us$ 时间内的 $v-t$ 图线为半圆，质点在该段时间内的平均速度为\xzanswer{B}
\onepicture[h=3.2cm, align=right]{\includesvg{figs/06}}
\fourchoices[answer=B]
{$\pi\Ums$}
{$\frac{\pi}{2}\Ums$}
{$1\Ums$}
{$\frac{\pi}{4}\Ums$}
%% answer: B
%% memo:
\memoanswer{
$v-t$ 图线与 $t$ 轴所围“面积”的物理意义即该段时间内物体运动的位移，$s=\frac{1}{2}S_{\text{圆}}=\frac{\pi r^{2}}{2}=2\pi$，平均速度为 $\overline{v}=\frac{s}{t}=\frac{\pi}{2}\Ums$。正确选项为 B。
}

\item
%% number: 7
%% paperName: 2021年上海高考第 7 题 3 分
%% typeId: 1
%% type: 单选题
%% chapter: 电场
%% point: 电势能电势差电场力做功
%% method: 无
%% score: 3
%% degree: 930
%% duplicateId: 0
%% body:
带正电的粒子在沿电场线方向运动的过程中，必定减小的物理量是\xzanswer{C}
\fourchoices[answer=C]
{粒子所受的电场力}
{粒子的速度}
{粒子具有的电势能}
{粒子的加速度}
%% answer: C
%% memo:
\memoanswer{
正电荷受到的电场力方向与电场方向相同，因此电场力做正功，电势能减小。正确选项为 C。
}

\item
%% number: 8
%% paperName: 2021年上海高考第 8 题 3 分
%% typeId: 1
%% type: 单选题
%% chapter: 气体性质
%% point: 气体图象
%% method: 无
%% score: 3
%% degree: 801
%% duplicateId: 0
%% body:
玻璃管开口向下竖直放置，管内用水银封闭一定质量的理想气体。在玻璃管绕顶端缓慢转到虚线所示位置的过程中，管内封闭气体状态变化可能是下图中的\xzanswer{D}
\onepicture[h=3.8cm, align=right]{\includesvg{figs/08}}
\fourchoices[answer=D, ispicture=true, h=2.6cm]
{\includesvg{figs/08a}}
{\includesvg{figs/08b}}
{\includesvg{figs/08c}}
{\includesvg{figs/08d}}
%% answer: D
%% memo:
\memoanswer{
玻璃管绕顶端\textbf{缓慢}转动，可以看成等温变化。转动过程中水银柱的竖直高度减小，产生的压强也减小，气体压强 $p=p_{0}-\rho gh$ 随之增大，由玻意耳定律可知，气体体积减小。正确选项为 D。
}

\item
%% number: 9
%% paperName: 2021年上海高考第 9 题 4 分
%% typeId: 1
%% type: 单选题
%% chapter: 直线运动
%% point: 匀变速直线运动
%% method: 无
%% score: 4
%% degree: 940
%% duplicateId: 0
%% body:
小球做直线运动的频闪照片如图所示，由此可以断定小球的\xzanswer{A}
\onepicture[w=6.5cm]{\includesvg{figs/09}}
\fourchoices[answer=A]
{加速度向右}
{速度向右}
{加速度向左}
{速度向左}
%% answer: A
%% memo:
\memoanswer{
若小球向右运动，由频闪照片可知单位时间内通过的位移越来越大，小球做加速运动，加速度方向向右；若小球向左运动，单位时间内通过的位移越来越小，小球做减速运动，加速度方向仍向右。正确选项为 A。
}

\item
%% number: 10
%% paperName: 2021年上海高考第 10 题 4 分
%% typeId: 1
%% type: 单选题
%% chapter: 磁场
%% point: 电流的磁场
%% method: 无
%% score: 4
%% degree: 666
%% duplicateId: 0
%% body:
长直导线中通以与时间成正比的电流，该电流会在其周围产生\xzanswer{D}
\fourchoices[answer=D]
{不随时间变化的匀强磁场}
{随时间变化的匀强磁场}
{不随时间变化的非匀强磁场}
{随时间变化的非匀强磁场}
%% answer: D
%% memo:
\memoanswer{
直线电流产生的磁感线为以直导线为圆心的一组不等距同心圆环，是非匀强磁场，该磁场的磁感应强度随直导线中电流的增大而增大。正确选项为 D。
}

\item
%% number: 11
%% paperName: 2021年上海高考第 11 题 4 分
%% typeId: 1
%% type: 单选题
%% chapter: 曲线运动
%% point: 竖直平面圆周运动
%% method: 无
%% score: 4
%% degree: 790
%% duplicateId: 0
%% body:
半径为 $R$，内壁光滑的圆轨道固定在竖直平面内，小球以 $\sqrt{2gR}$（$g$ 为重力加速度）的速度经过圆轨道最低点后到达轨道上与圆心等高处时的速度为 $v$，加速度为 $a$，则\xzanswer{C}
\fourchoices[answer=C]
{$v=0$，$a=0$}
{$v\neq0$，$a=0$}
{$v=0$，$a\neq0$}
{$v\neq0$，$a\neq0$}
%% answer: C
%% memo:
\memoanswer{
由机械能守恒定律：$\frac{1}{2}m(\sqrt{2gR})^{2}=mgR+\frac{1}{2}mv^{2}$，解得 $v=0$；由于 $v=0$，所以 $F_{\text{向}}=0$，即此时内壁对小球没有弹力作用，小球只受重力，$a=g\neq0$。正确选项为 C。
}

\item
%% number: 12
%% paperName: 2021年上海高考第 12 题 4 分
%% typeId: 1
%% type: 单选题
%% chapter: 电场
%% point: 电场强度
%% method: 无
%% score: 4
%% degree: 658
%% duplicateId: 0
%% body:
如图，在点电荷 $Q$ 附近不同位置放置检验电荷 a、b，描述 a、b 所带电量 $q$ 与所受电场力大小 $F$ 的两组数据可能是下图中的\xzanswer{B}
\onepicture[h=1.4cm, align=right]{\includesvg{figs/12}}
\fourchoices[answer=B, ispicture=true, h=2.6cm]
{\includesvg{figs/12a}}
{\includesvg{figs/12b}}
{\includesvg{figs/12c}}
{\includesvg{figs/12d}}
%% answer: B
%% memo:
\memoanswer{
\onepicture[h=2.6cm, align=right]{\includesvg{figs/12_an}}

由 $E=\frac{F}{q}=\frac{kQ}{r^{2}}$ 可知，$F-q$ 图中的点的斜率表示该点所在位置处的电场强度，且距离场源电荷 $Q$ 越远，斜率越小。正确选项为 B。
}

\item
%% number: 13
%% paperName: 2021年上海高考第 13 题 4 分
%% typeId: 3
%% type: 填空题
%% chapter: 原子和原子核
%% point: 半衰期
%% method: 无
%% score: 4
%% degree: 920
%% duplicateId: 0
%% body:
氚核（\ce{^{3}_{1}H}）的中子数为\tkanswer[1.2cm]{2}，半衰期为 12 年，一定量的氚核经过\tkanswer[1.2cm]{24}年减少到原来的 $1/4$。
%% answer: 2，24
%% memo:
\memoanswer{
（1）中子数为核子数减去质子数，即 $3-1=2$ 个；

（2）根据半衰期的定义，氚核减少到原来的 $1/4$ 用时 2 个半衰期，即 24 年。
}

\item
%% number: 14
%% paperName: 2021年上海高考第 14 题 4 分
%% typeId: 3
%% type: 填空题
%% chapter: 气体性质
%% point: 查理定律
%% method: 无
%% score: 4
%% degree: 840
%% duplicateId: 0
%% body:
密闭钢瓶内温度降低，瓶内气体分子的\tkanswer[1.8cm]{平均动能}减小，对瓶壁单位面积的平均作用力\tkanswer[1.2cm]{减小}（选填：“增加”、“不变”或“减小”）。
%% answer: 平均动能，减小
%% memo:
\memoanswer{
（1）温度是分子平均动能的标志，温度降低，平均动能减小；

（2）由查理定律可知，在体积不变的情况下，一定质量气体的压强随温度的降低而减小，即对瓶壁单位面积的平均作用力减小。
}

\item
%% number: 15
%% paperName: 2021年上海高考第 15 题 4 分
%% typeId: 3
%% type: 填空题
%% chapter: 光学
%% point: 光的电磁说
%% method: 无
%% score: 4
%% degree: 810
%% duplicateId: 0
%% body:
上海科学家研发出的激光器能在 $2\times10^{-14}\Us$ 内输出能量为 $400\UJ$ 的超强超短激光，这段时间内激光器的平均功率高达\tkanswer[3.4cm]{$2\times10^{16}$} $\UW$；该激光的中心波长为 $8\times10^{-7}\Um$，对应的频率为\tkanswer[3.4cm]{$3.75\times10^{14}$} $\UHz$。
%% answer: 2×10^16，3.75×10^14
%% memo:
\memoanswer{
（1）$P=\frac{W}{t}=\frac{400}{2\times10^{-14}}\UW=2\times10^{16}\UW$；

（2）$f=\frac{c}{\lambda}=\frac{3\times10^{8}}{8\times10^{-7}}\UHz=3.75\times10^{14}\UHz$。
}

\item
%% number: 16
%% paperName: 2021年上海高考第 16 题 4 分
%% typeId: 3
%% type: 填空题
%% chapter: 运动和力
%% point: 力学单位制
%% method: 量纲分析
%% score: 4
%% degree: 630
%% duplicateId: 0
%% body:
飞行器在太空航行时，发动机的喷气速率 $v$ 在一定的近似条件下仅由下列四个量决定：高温气体的温度 $T$（单位：$\UK$）、摩尔质量 $M$（单位：$\Ukgmol$）、定压摩尔热容 $C$（单位：$\UJKmol$）、无单位的比例系数 $\alpha$。若使用国际单位制中的基本单位，$C$ 的单位又可表示为\tkanswer[4.6cm]{$\Ukgmqskmol$}；用 $T$、$M$、$C$ 和 $\alpha$ 表示的 $v=\alpha\left(\frac{CT}{M}\right)^{\lambda}$，式中 $\lambda=$\tkanswer[1.0cm]{0.5}。
%% answer: kg·m^2·s^-2·K^-1·mol^-1，0.5
%% memo:
\memoanswer{
（1）由 $W=Fs$ 和 $F=ma$ 可知，单位焦耳（J）可以表示为 $\UNcdotm$，并进一步表示为 $\mathrm{kg}\cdot\mathrm{m}^{2}\cdot\mathrm{s}^{-2}$，因此 $C$ 的单位 $\UJKmol$ 可以表示为 $\Ukgmqskmol$；

（2）$\alpha\left(\frac{CT}{M}\right)^{\lambda}$ 的单位可以表示为 $[\mathrm{kg}\cdot\mathrm{m}^{2}\cdot\mathrm{s}^{-2}\cdot\mathrm{K}^{-1}\cdot\mathrm{mol}^{-1}][\mathrm{K}][\mathrm{kg}^{-1}\cdot\mathrm{mol}]=[\mathrm{m}^{2}\cdot\mathrm{s}^{-2}]$，而速度 $v$ 的单位可以表示为 $[\mathrm{m}\cdot\mathrm{s}^{-1}]$，比较两者可知 $\lambda=0.5$。
}

\item
%% number: 17
%% paperName: 2021年上海高考第 17 题 4 分
%% typeId: 3
%% type: 填空题
%% chapter: 机械振动
%% point: 单摆
%% method: 无
%% score: 4
%% degree: 540
%% duplicateId: 0
%% body:
如图为某测量引力常量 $G$ 的实验原理示意图，质量为 $M$ 的均匀球 A 固定，可视为质点的小球 B 由轻质细线悬挂于 O 点。平衡时两球心等高、间距为 $r$，BO 间水平距离为 $d$（远小于细线长度）。迅速移去 A 球后，B 球的摆动\tkanswer[1.2cm]{可以}（选填：“可以”或“不可以”）视为简谐运动。测量 $M$、$r$、$d$ 和 B 球的摆动周期 $T$ 后可得 $G=$\tkanswer[3.2cm]{$\frac{4\pi^{2}r^{2}d}{MT^{2}}$}。（$\theta$ 很小时，$\sin\theta\approx\tan\theta$）
\onepicture[h=3.6cm, align=right]{\includesvg{figs/17}}
%% answer: 可以，4π²r²d/(MT²)
%% memo:
\memoanswer{
（1）由于 BO 间水平距离为 $d$ 远小于细线长度，因此摆角 $\theta$ 也很小，可以看成小于 $5^\circ$，B 球的摆动可以视为简谐运动；

（2）开始时 B 球在三个力的作用下平衡，有：$G\frac{Mm}{r^{2}}=mg\tan\theta$，可求得 $G=\frac{gr^{2}\tan\theta}{M}\approx\frac{gr^{2}\sin\theta}{M}=\frac{gr^{2}d}{Ml}$；

移去 A 球后，B 球做简谐振动，周期 $T=2\pi\sqrt{\frac{l}{g}}$，可求得：$g=\frac{4\pi^{2}l}{T^{2}}$，代入上式可得：$G=\frac{4\pi^{2}r^{2}d}{MT^{2}}$。
}

\item
%% number: 18
%% paperName: 2021年上海高考第 18 题 10 分
%% typeId: 4
%% type: 实验题
%% chapter: 电场
%% point: 描绘等势面
%% method: 无
%% score: 10
%% degree: 651
%% duplicateId: 0
%% body:
“用 DIS 描绘电场的等势线”的实验示意图如图所示。图中 A、B 是连接电源的两个电极，基准点 c 位于 A、B 连线的中点，f、d 连线和 A、B 连线垂直。
\begin{enumerate}
	\item
	实验时使用的是\tkanswer[1.5cm]{电压}传感器。
	\item
	若传感器的红、黑色探针分别接触图中的 d、f 两点时，传感器示数小于零；保持红色探针位置不变，将黑色探针从 f 点换到 e 点，传感器示数\tkanswer[1.2cm]{大于}（选填：“大于”或“小于”）零。
	\item
	在图上画出过 f 点的等势线。
	\item
	（单选题）实验中得到的曲线实际上是\tkanswer[1.0cm]{C}。

	\fourchoices[answer=C]
	{一对等量同种电荷的静电场的等势线}
	{一对等量异种电荷的静电场的等势线}
	{稳恒电流场中的等势线}
	{变化电流场中的等势线}
\end{enumerate}
\onepicture[h=3.8cm, align=right]{\includesvg{figs/18}}
%% answer: （1）电压（2）大于（3）见详解（4）C
\jdanswer{
\begin{enumerate}
	\item
	电压

	\item
	大于

	\item
	如图所示

	\item
	C
\end{enumerate}
}
%% memo:
\memoanswer{
（1）本实验测量的是两点间的电势差，即电压，因此使用的是电压传感器；

（2）“传感器的红、黑色探针分别接触图中的 d、f 两点时，传感器示数小于零”，说明红色探针所在位置 d 的电势较低，可以推断 A 电极与电源正极相连，进一步可知 $\varphi_{d}>\varphi_{e}$，将黑色探针换到 e 点，测量的是低电势，所以示数大于零；

（3）如图所示：

\onepicture[h=3.0cm]{\includesvg{figs/18_an}}

（4）这个实验通过电源接上接线柱，在导电物质上产生的稳恒电流场模拟一对等量异种点电荷产生的静电场。正确选项为 C。
}

\item
%% number: 19
%% paperName: 2021年上海高考第 19 题 15 分
%% typeId: 6
%% type: 计算题
%% chapter: 运动和力
%% point: 牛顿定律的应用
%% method: 无
%% score: 15
%% degree: 770
%% duplicateId: 0
%% body:
如图，倾角为 $\theta$ 的足够长斜面 ABC 上，AB 段光滑，长为 $l$，BC 段粗糙。一质量为 $m$ 的物体（视为质点）在平行于斜面的拉力作用下自 A 由静止开始做匀加速直线运动，经时间 $t_{0}$ 到达 B。（重力加速度为 $g$，BC 段与物体间的动摩擦因数为定值）
\begin{enumerate}
	\item
	求拉力的大小 $F$；
	\item
	物体到 B 点时拉力的功率 $P_{B}$ 为多大？
	\item
	从到达 B 点开始，拉力的功率保持为 $P_{B}$，物体做减速运动。定性画出物体由 A 运动到 C 过程的 $v-t$ 图像。（仅要求作出图像）
\end{enumerate}
\onepicture[w=7.5cm, align=right]{\includesvg{figs/19}}
%% answer: 
\jdanswer{
\begin{enumerate}
	\item
	$F=mg\sin\theta+\frac{2ml}{t_{0}^{2}}$

	\item
	$P_{B}=\frac{2mgl\sin\theta}{t_{0}}+\frac{4ml^{2}}{t_{0}^{3}}$

	\item
	如图所示（见详解）
\end{enumerate}
}
%% memo:
\memoanswer{
（1）受力分析如图所示，由牛顿第二定律可得：

$F-mg\sin\theta=ma$

$a=\frac{F}{m}-g\sin\theta$

由匀变速直线运动公式可得：

$l=\frac{1}{2}at_{0}^{2}=\frac{1}{2}\left(\frac{F}{m}-g\sin\theta\right)t_{0}^{2}$

解得：$F=mg\sin\theta+\frac{2ml}{t_{0}^{2}}$

\onepicture[h=2.6cm]{\includesvg{figs/19_an1}}

（2）由匀变速直线运动公式可得：$l=\frac{v_{B}}{2}t_{0}$，解得：$v_{B}=\frac{2l}{t_{0}}$

拉力功率：$P_{B}=Fv_{B}=\frac{2mgl\sin\theta}{t_{0}}+\frac{4ml^{2}}{t_{0}^{3}}$

（3）如图所示

\onepicture[h=2.4cm]{\includesvg{figs/19_an2}}

【分析过程】AB 过程中，物体做初速为零的匀加速直线运动，图线为过原点的倾斜直线。

BC 过程中，由 $F=\frac{P}{v}$ 可知，在 $P$ 不变的情况下，$v$ 减小，导致 $F$ 增大；

由牛顿第二定律：$mg\sin\theta+f-F=ma$，$F$ 增大，合外力减小，加速度减小，物体做加速度不断减小的减速运动，图线为斜率减小的曲线。

理论上需要无穷长的时间才达到匀速，图线末端最好画渐近线。
}

\item
%% number: 20
%% paperName: 2021年上海高考第 20 题 15 分
%% typeId: 6
%% type: 计算题
%% chapter: 电磁感应
%% point: 电磁感应中的变加速
%% method: 无
%% score: 15
%% degree: 540
%% duplicateId: 0
%% body:
两竖直固定的光滑金属导轨间距 $L=0.75\Um$，磁感应强度 $B=0.8\UT$ 的匀强磁场与两导轨所在平面垂直，一金属棒与导轨垂直放置。导轨上方接有如图（a）所示电路，其中 R 为一电阻性元件，定值电阻阻值 $R_{0}=10\UO$。断开电键 S，金属棒从静止开始下滑直到匀速运动，此期间由电压、电流传感器得到 R 的 $U-I$ 图线如图（b）所示。（导轨足够长，导轨和金属棒的电阻均忽略不计，金属棒下滑过程中始终与导轨接触良好，重力加速度 $g$ 取 $10\Umsq$）
\begin{enumerate}
	\item
	求金属棒达到匀速运动时的速度 $v$；
	\item
	求金属棒的质量 $m$；
	\item
	在金属棒匀速运动时合上 S，求经过足够长时间后金属棒的动能 $E_{k}$。
\end{enumerate}
\onepicture[w=9.5cm, align=right]{\includesvg{figs/20}}
%% answer: 
\jdanswer{
\begin{enumerate}
	\item
	$v=10\Ums$

	\item
	$m=0.03\Ukg$

	\item
	$E_{k}=\frac{1}{6}\UJ$
\end{enumerate}
}
%% memo:
\memoanswer{
（1）由图（b）可知，当金属棒达到匀速时，R 两端的电压稳定在 $6\UV$，且导轨和金属棒的电阻忽略不计。则有：

$E=U=BLv$

$v=\frac{U}{BL}=\frac{6}{0.8\times0.75}\Ums=10\Ums$

（2）金属棒匀速运动时处于受力平衡状态，由图（b）可知此时通过金属棒的电流 $I=0.5\UA$。可得：

$mg=BIL$

$m=\frac{BIL}{g}=\frac{0.8\times0.5\times0.75}{10}\Ukg=0.03\Ukg$

（3）（合上 S 后，金属棒做加速度不断减小的减速运动）经过足够长的时间后，棒仍做匀速直线运动，此时仍处于受力平衡状态，因此通过棒的电流不变，仍为 $I=0.5\UA$。

设此时电阻 R 两端的电压为 $U_{R}$，通过电阻 R 的电流为 $I_{R}$，通过 $R_{0}$ 的电流为 $I_{0}$，应满足：

$I_{R}=I-I_{0}$

$I_{R}=0.5-\frac{U_{R}}{R_{0}}$

$U_{R}=5-10I_{R}$

\onepicture[h=3.0cm]{\includesvg{figs/20_an}}

在图（b）中画出这个方程对应的图线，与 R 的伏安特性曲线的交点的意义就是此时的工作点，解得 $U_{R}=2\UV$，$I_{R}=0.3\UA$。

设此时匀速运动的速度大小为 $v'$，则有：

$U_{R}=BLv'$

$v'=\frac{U_{R}}{BL}=\frac{2}{0.8\times0.75}\Ums=\frac{10}{3}\Ums$

$E_{k}=\frac{1}{2}mv'^{2}=0.5\times0.03\times\left(\frac{10}{3}\right)^{2}\UJ=\frac{1}{6}\UJ$
}

\end{enumerate}

\end{document}
